\chapter{Compositional Specifications for Version Compatibility}\label{ch:compositional}

\section{From a local pass to a compositional one}
\label{sec:comp-motivation}

The inference instrument of \cref{ch:inferrer} reasons about one method at a
time. Its heuristics read the syntax of a single method body---a parameter
dereferenced without a guard, an array indexed by a loop variable, an early
\texttt{throw} that validates an argument---and emit a contract for that method
alone. This locality is what makes the instrument cheap and deterministic, but
it is also its ceiling. A method's true precondition is rarely a property of its
own body only; it is inherited from the methods it calls. When a caller invokes
a callee whose contract requires something of its arguments, that requirement
becomes, after substituting the actual arguments, an obligation on the caller.
Recovering such inherited obligations is the business of \emph{compositional},
rather than local, inference.

The inference instrument already performs a limited form of this cross-method
reasoning. During its single local pass it consults an
\texttt{InterproceduralAnalyzer}: when it infers a method \texttt{m} that
contains a call \texttt{n(args)}, it looks up \texttt{n}'s cached contract,
substitutes the formal parameters of \texttt{n} with the actual arguments at the
call site, and appends the substituted preconditions to \texttt{m}'s set. This
is \emph{single-pass propagation}: each method is visited once, each callee is
read once, and its contract is pulled into the caller atomically. It is sound
for the simple case---an obligation the callee requires \emph{unconditionally}
must hold wherever the callee is called---and it is fast and order-independent
provided callees are inferred before callers.

The question this chapter answers, and the question that anchors the thesis's
second contribution, is whether that single pass captures everything a genuinely
compositional analysis would, or whether a second, deliberately compositional
pass recovers obligations the first structurally cannot express. The answer
matters twice over. If a second pass adds nothing, the simpler architecture is
vindicated and the engineering cost of a second pass is unjustified. If it adds
a great deal, then the instrument of \cref{ch:inferrer} was leaving recoverable
specifications on the table, and---more consequentially---the same machinery
that lifts a callee's contract into its caller is exactly what lets a
\emph{change} to a callee's contract be propagated across a call graph, which is
the mechanism on which behavioural version-compatibility reasoning rests
(\cref{sec:comp-versioning}).

\section{The weakest-precondition intuition}
\label{sec:comp-wp}

Where the single pass is blind is at call sites that are not reached
unconditionally. Dijkstra's weakest-precondition calculus
\citep{dijkstra1976discipline}, resting on the axiomatic foundations of Hoare
\citep{hoare1969axiomatic}, gives the correct account: to compute what must hold
at a method's entry, one walks the body backwards, accumulating the guard
context that governs each program point. A statement reached only when a
condition \texttt{g} holds contributes not its raw obligation \texttt{p} but the
implication \texttt{g \(\Rightarrow\) p}. A small example makes the gap visible.
Consider a callee with an inferred, verified precondition and a caller that
invokes it on one branch only:

\begin{lstlisting}[language=Java,basicstyle=\ttfamily\small,frame=none]
//@ requires divisor != 0;
static int safeDivide(int a, int divisor) { return a / divisor; }

static int scale(int value, int factor) {
    if (factor != 0) {
        return safeDivide(value, factor);
    }
    return value;
}
\end{lstlisting}

At the call site \texttt{safeDivide(value, factor)}, substitution yields the
clause \texttt{factor != 0}. The single pass now faces a dilemma with no correct
resolution. Conjoining \texttt{factor != 0} unconditionally into \texttt{scale}'s
precondition would forbid the very call \texttt{scale(value, 0)} that the method
explicitly and correctly handles on its other branch; skipping the clause to
avoid that error loses the obligation entirely on the path where the division
does fire. A compositional walk resolves the dilemma by lifting the substituted
clause through the enclosing guard, producing the branch-conditional clause
\texttt{factor != 0 \(\Rightarrow\) factor != 0}. Here the guard happens to
entail the obligation, so the lifted clause is a tautology and the example is
deliberately trivial; in real code the guard and the callee's requirement
differ, and a call to a method requiring a non-empty list, guarded by a size
check, lifts to \texttt{size > 0 \(\Rightarrow\) list.size() > 0}---a clause
that carries genuine information the single pass cannot state in any form.

A second structural class is equally out of reach. When a caller invokes
\texttt{x.n()} on a receiver whose declared type has several implementations of
\texttt{n}, runtime dispatch may select any of them, each with its own inferred
precondition. Single-pass propagation resolves the call to one nominal callee
and sees only that one. A compositional pass, aware of the class hierarchy,
enumerates the candidates and emits the disjunction
\texttt{(p\(_1\) || \dots\ || p\(_k\))} of their substituted preconditions.
This is a deliberate \emph{weakening} of the caller's obligation rather than a
sound safety guarantee---the heuristically inferred per-override contracts need
not respect behavioural subtyping \citep{mostafa2017behavioral}, so the strictly
sound conjunction is frequently spuriously strong---but it is a shape a single
nominal resolution cannot produce at all.

\section{The compositional refinement pass}
\label{sec:comp-pass}

The realised design is a lifting step layered over the existing local result,
not a full weakest-precondition transformer. A complete transformer would
require a decision procedure and a precise operational semantics for every Java
construct, which is the province of full verifiers such as OpenJML's static
checker \citep{cok2014openjml,leavens2006jml}, not of a heuristic inferrer. The
pass instead takes as input the populated specification cache from
\cref{ch:inferrer}, a call graph over the same methods, and the abstract syntax
tree of each declaration, and it mutates the cache in place.

\begin{figure}[t]
\centering
\begin{tikzpicture}[
  font=\footnotesize\sffamily,
  >={Stealth[length=1.6mm,width=1.4mm]},
  node distance=7mm and 12mm,
  box/.style={rectangle, rounded corners=1.5pt, draw=black!60,
    minimum height=6mm, minimum width=16mm, inner sep=3pt, align=center},
]
  \node[box] (leaf) {leaf callee\\(refined first)};
  \node[box, right=of leaf] (mid) {intermediate\\caller};
  \node[box, right=of mid] (root) {entry-point\\caller};
  \draw[->] (mid) -- node[above,font=\tiny]{calls} (leaf);
  \draw[->] (root) -- node[above,font=\tiny]{calls} (mid);
  \draw[->, dashed, bend left=32] (leaf.north) to
    node[above,font=\tiny]{contract lifts to} (mid.north);
  \draw[->, dashed, bend left=32] (mid.north) to
    node[above,font=\tiny]{contract lifts to} (root.north);
\end{tikzpicture}
\caption{Refinement direction. The driver visits strongly-connected components
of the call graph in reverse topological order, so each callee is refined before
any of its callers; a callee's refined contract then lifts into the caller as a
substituted, guard-conditioned obligation.}
\label{fig:comp-callgraph}
\end{figure}

A driver iterates the strongly-connected components of the call graph in reverse
topological order, so that callees are refined before their callers
(\cref{fig:comp-callgraph}). A component with no recursion is refined in a single
sweep; a mutually recursive component is refined repeatedly until no method gains
a new clause, or until a configurable cap (sixteen iterations) bounds the runtime
on pathological graphs. The two-tier treatment is forced by the graph rather than
chosen for convenience: on an acyclic graph one reverse-topological sweep already
refines every method against fully refined callees, but a cycle admits no such
order, so a method first visited against a still-stale callee must be revisited
once that callee is itself refined. The iteration terminates because each
method's clause set only grows and is bounded above by the finitely many
substituted callee clauses reachable in the component.

For each method the refiner walks every call expression in the body and, at each
call site, resolves the callee, substitutes the actual arguments for the callee's
formal parameters, lifts the substituted clause through the conjunction of any
enclosing \texttt{if} or ternary guards (with \texttt{else} branches
contributing the negated guard), detects polymorphic dispatch and emits a
disjunction over candidates where it occurs, and conjoins the result into the
caller's set if it is not already present. Substitution is textual with word
boundaries, so a formal parameter \texttt{s} does not clobber a substring of
\texttt{this.size}. Two soundness gates reject unsafe substitutions outright:
a call whose arguments or whose enclosing guards are side-effecting (an
increment, a decrement) is skipped, because substituting an expression whose
evaluation changes state does not preserve meaning. Postconditions are propagated
by the same machinery but gated on a termination flag---a non-terminating
callee's postcondition describes a state established only if the callee returns,
so propagating it unconditionally would be unsound.

\section{Entry-state well-formedness}
\label{sec:comp-wellformed}

A lifted clause is only admissible if it can be \emph{stated} at the caller's
entry. A JML \texttt{requires} clause may reference only entry-state names---the
method's parameters and the fields visible to it---and never a body-local such as
a loop index or a temporary, because a precondition is evaluated before the body
runs and body-locals do not yet exist. Two routes threaten this discipline.
Lifting a guard that tests a body-local (\texttt{if (i < n) \dots}, where
\texttt{i} is a loop variable) would drag \texttt{i} into a \texttt{requires}
clause; and substituting a caller-local argument into a callee's precondition
would do the same by the other route. Either produces a clause that is
syntactically ill-formed as a precondition, no matter how faithful it is to the
program's meaning.

The pass therefore checks every candidate clause against the caller's entry scope
before emitting it, discarding any clause that mentions a name not in scope at
entry. This filter---realised as the \texttt{EntryScopeFilter} component---is not
cosmetic. Enforcing it materially lowered the recovered count: an earlier,
unfiltered version of the analysis reported on the order of forty-five thousand
new preconditions, but many of those referenced body-locals and were not
legitimate entry-state obligations at all. With the filter enforced, the pass
yields roughly \textbf{18{,}854 additional well-formed preconditions} across the
study corpus (\cref{sec:comp-measurement}), and a post-hoc re-scan confirms that
none of the emitted clauses references a body-local on any library. The corrected
figure is the one this thesis reports and defends.

The distinction that governs the whole chapter must be stated plainly here. A
\emph{well-formed} clause is syntactically valid JML and scope-correct: it names
only entry-state values and could, in principle, be handed to a verifier. A
\emph{verified} clause is one an SMT-backed checker has actually discharged
\citep{cok2014openjml,demoura2008z3}. The compositional pass produces the former,
not the latter. It is a structural transformation over already-cached contracts:
given the callee specifications, whatever their semantic quality, it produces new
caller clauses and confirms they are expressible at entry. It does not run
OpenJML on them, and it makes no claim that the 18{,}854 clauses are individually
sound against the code's intent---that is the inference instrument's
responsibility, discharged by the verification discipline of \cref{ch:inferrer},
not the compositional pass's. Throughout this chapter, therefore, the headline
number is a count of \emph{well-formed} preconditions, not verified ones, and the
two are never conflated.

\section{Measurement study}
\label{sec:comp-measurement}

The pass was measured, not proved. The study asks how many additional well-formed
preconditions the second pass recovers, over what fraction of the call graph, on
real Java. The corpus is the source of five open-source libraries chosen to span
four idiom families: Apache Commons Lang and Commons IO (imperative
utility-class and I/O idioms), Commons Math (numerical algorithms), and jOOL and
Vavr (higher-order, generics-heavy functional idioms). Commons Lang and Commons
IO are shared with the studies of \cref{ch:inferrer}, which keeps cross-chapter
comparisons free of corpus-difference confounds. Combined, the corpus is 1{,}712
source files and 21{,}052 methods that carry a non-empty local baseline contract.
For each library the harness parses every source file, builds the call graph,
runs the local inference (including its single-pass interprocedural propagation)
to snapshot the baseline, invokes the compositional pass, and diffs each method's
refined contract against its baseline, counting and categorising the newly added
clauses. Every measurement is deterministic given the inputs and is reported
once.

The second pass strictly extends the baseline on every library
(\cref{tab:comp-headline}). It adds at least one clause to between roughly 21 and
39 per cent of cached methods---median about a third---and in total recovers
\textbf{18{,}854 new, well-formed precondition clauses} (together with a further
body of new postcondition clauses) across the five libraries. The effect is a
population-level one rather than the artefact of a few enormous methods: the mean
number of new preconditions per refined method is modest, in the low single
digits on every library, well below the per-library maxima, which are therefore
deep-tail outliers. Crucially, this gap persists even though every library's
baseline already invokes single-pass interprocedural propagation; the second pass
discovers obligations the first cannot express, not obligations the first simply
failed to compute.

\begin{table}[t]
\centering
\small
\caption{Effect of the compositional pass over the local baseline, per library.
``Refined'' methods gained at least one new clause; every emitted clause is
well-formed (entry-state-expressible), not separately verified.}
\label{tab:comp-headline}
\begin{tabular}{lrrrrr}
\toprule
& \textbf{Lang} & \textbf{IO} & \textbf{Math} & \textbf{jOOL} & \textbf{Vavr} \\
\midrule
Methods in cache        & 3{,}494 & 2{,}111 & 6{,}645 & 3{,}691 & 5{,}111 \\
Methods refined         & 1{,}367 & 711     & 2{,}203 & 784     & 1{,}541 \\
\hspace{1em}\% of cache & 39.1    & 33.7    & 33.2    & 21.2    & 30.2 \\
New precondition clauses & 2{,}632 & 1{,}329 & 8{,}966 & 4{,}262 & 1{,}665 \\
\bottomrule
\end{tabular}
\end{table}

The shape of the additions explains \emph{why} the single pass falls short.
Categorising each new precondition by its syntactic form, the two shapes that lie
outside the single pass's reach by construction---branch-conditional implications
(\texttt{guard \(\Rightarrow\) obligation}) and disjunctions from polymorphic
dispatch---are present on every library and together account for between two-fifths
and three-fifths of new clauses on each. Branch-conditional clauses are the
single largest shape on three of the five libraries, exceeding two-fifths on
Commons Lang and on Vavr, and remain at least a third on the other two, where
pervasive input validation instead pushes plain null checks to the plurality (the
matrix and array arguments of Commons Math, the lambda pinch-points of jOOL).
Disjunctive clauses, the polymorphic-dispatch signal, peak on Vavr, the library
with the deepest interface-inheritance chains. Null and range obligations also
appear, but these are shapes the single pass \emph{would} have produced at an
unguarded call site; the second pass adds them only because the call site was
guarded by a condition the single pass would not lift. A syntactic vacuity test
confirms that the branch-conditional counts are not inflated by tautologies of
the \texttt{g \(\Rightarrow\) g} form: such vacuous clauses are a negligible
fraction of the branch-conditional total on every library.

The pass is cheap. Across all five libraries the refinement runs in around fifty
seconds in total---on the order of a couple of milliseconds per method---an
addition of seconds to a local inference that already takes minutes on libraries
of this size. The fixpoint iteration never fired: under the analyser's
signature-matching scheme all five call graphs are acyclic, so each method is
refined exactly once and the sixteen-iteration cap is never approached. This is a
statement about precondition propagation along the call graph specifically, and
it does not generalise to every specification a compositional inferrer might
produce; class invariants, whose inference is a greatest-fixpoint driven by
shared-state coupling rather than by call edges, are a genuinely harder problem
left to future work and revisited in \cref{ch:discussion}. Two caveats bound the
headline. Callee resolution uses simple-name matching, which rejects some
overloaded call sites on arity mismatch, so unmatched sites can only ever add
clauses the count omits; the reported figure is thus a \emph{lower bound} on the
true compositional gain. And clause equality is syntactic, so the count is an
upper bound on the number of distinct \emph{logical} obligations. Both caveats
concern the count's precision, not the qualitative finding, which is robust: the
single pass is not subsuming on any library measured.

\section{Behavioural version compatibility}
\label{sec:comp-versioning}

The mechanism that recovers a caller's obligations from a callee's contract is,
turned around, a mechanism for propagating a \emph{change} in a callee's contract
to everything that depends on it. This is the pivot from measurement to
application, and it is the concern that led the thesis from the outset: library
and version compatibility.

Modern software is assembled from libraries, and the contract between a library
and its clients is signalled by a version number under the convention of semantic
versioning \citep{preston2013semver}, which promises that a compatible release
will not break existing callers. The convention is honoured unevenly: breaking
changes are routinely shipped in releases whose version number declares them
compatible \citep{raemaekers2017semver}, precisely because the contract being
versioned is implicit. There is no machine-checkable statement of what a release
guarantees, and hence no mechanical way to decide whether a new release still
honours what the old one promised. The consequence falls hardest on exactly the
systems whose behaviour is least understood---the legacy systems that most matter
\citep{feathers2004legacy}---and it is amplified by any lifecycle that regenerates
components against informal descriptions, where every regeneration is a potential
silent breaking change.

Inferred, compositional contracts turn this asserted convention into a checkable
property. Given two consecutive releases of a library, the instrument infers a
contract for each method in each release; diffing the contracts method by method
yields the set of methods whose behavioural interface has changed, and the
compositional pass then propagates each change across the call graph to compute
its \emph{blast radius}---the set of callers whose own obligations shift as a
consequence. A precondition that a release has \emph{strengthened} is the sharp
case: client code that satisfied the old requirement may now violate the new one,
which is a breaking change in behaviour whether or not the version number admits
it. The analysis is entirely static and run-free---no client is executed, no test
is run---and so it applies to a release exactly as it sits in a repository.

A concrete version step demonstrates the mechanism operating at scale. Inferring
contracts across two consecutive Commons Lang releases separated by roughly two
years, and diffing them, isolates a tractable candidate-breaking set: of the
3{,}188 methods present in both releases, around 525 (some 16.5 per cent) have an
inferred contract that differs---exactly the set a behavioural-compatibility
checker should surface. Of these, on the order of 144 have a \emph{strengthened}
precondition, the class of change most likely to break an existing caller; a spot
check finds a mixture of real precondition tightenings, such as added non-null
guards on the instance fields of several classes, together with some inferrer
noise. Two properties make this output useful even before each candidate is
validated against the release notes. It is the right \emph{size}---hundreds of
candidates out of thousands of methods, small enough for a human or a
downstream tool to triage---and it is the right \emph{shape}, classified into
strengthened, weakened, and mixed changes on both the precondition and
postcondition sides, so the class of each change is legible.

The reasoning is behavioural rather than syntactic, and this is the substance of
the contribution. A textual diff of two releases reports that a line changed; it
cannot say whether the change tightens what a caller must guarantee. A contract
diff, propagated compositionally, reports that a caller's obligation shifted and
in which direction---the semantics of the change, not its surface. This is what
it means to turn semantic versioning from a convention a maintainer asserts into
a property a tool can check, and it is the mechanism on which the AI-native
verification pipeline of \cref{ch:integration} relies to keep regenerated
components compatible with the callers that depend upon them.

The claim must be stated at its true strength, and no stronger. What the study
establishes is that the compositional pass recovers a large body of additional
\emph{well-formed}---not separately verified---preconditions, and that the same
propagation machinery yields a candidate-breaking set of the right size and shape
to make behavioural-compatibility triage tractable on a real version step. It
does not establish that every candidate is a true breaking change, nor that the
recovered clauses are individually verified; validating the candidate set against
release notes, and discharging the recovered clauses through OpenJML at corpus
scale, are the natural next steps set out in \cref{ch:conclusion}. The
contribution is a measured extension of the inference instrument and a working
static analysis for behavioural compatibility, presented as evidence rather than
as a proof of correctness.
